\section{Two accuracies, delayed projections, and the base reduction}
\label{sec:base}

The aim of this section is the active-cell minimum in
Equation~\eqref{eq:src-11}.  We first show that two fixed structured arguments
give an absolutely summable scale approximation, and use its finite
multipliers as predictor columns.  Their moments justify evaluating the
other projections on a finer partition.  Staggered approximation
accuracies then make every full pair small; the delayed projections
split into current projections and rolling differences with
square-energy bounds.  Estimating the term with two rolling differences
and removing inactive scales leaves the stated minimum.

Throughout this section, the four argument positions are distinct elements of
$\{0,1,2,3\}$.  We work on finite dyadic trees; an inherited projection family
is extended past the last decision by restriction and local refitting of its
existing columns.  This convention is useful when a projection evaluated on a
descendant partition occurs in a form associated with an ancestor interval.
No new column is required for such an extension.  All functions carrying
private coordinates are interpreted in the augmented Hilbert space from
Section~\ref{sec:compression}; the form $H_I$ uses only their physical
coordinates.

Here and below, a polynomial has an absolute degree.  Its coefficients may
depend on the fixed $q>2$ and on the fixed kernel class.  In particular, a
bound $\exp(k^{O(1)})$ permits these dependencies in its constant and in the
lower threshold for $k$, but its exponent is chosen independently of $q$.

\subsection{Absolute scale sums with two fixed structured arguments}

We first formulate precisely the class of inputs allowed in the two
remaining positions.  Fix a dyadic interval $J$ and $F\in L^2(J;\mathcal H)$.
An admissible sequence on a collection of descendants of $J$ is one of the
following:
\begin{enumerate}
\item the fixed input $F$;
\item $P_I F$, where $P_I$ belongs to a nested, cell-local column projection
family with path count, column size, and inverse private weight bounded by
$D_{\mathrm h}$;
\item a sequence $Z_I$, supported on $I$, for which
\[
 \sum_{I\subseteq J}|I|\|Z_I\|_{2,I}^2
 \le L^2\|\1_JF\|_2^2.
\]
\end{enumerate}
A sum of a fixed number of such sequences is also admissible.  The number
of summands is absorbed into absolute constants.  Here the index $I$ may
range over an arbitrary subcollection; projections on the missing cells
are completed by inheritance.  The estimates below remain valid when
the projection lists are enlarged: the additional dependence is through
$\log(2D_{\mathrm h})$, rather than a power of $D_{\mathrm h}$.

Fix two atoms of Lemma~\ref{lem:ap-separation}, in positions $i$ and $j$,
on their common birth subtree $J$.  Let $a,b$ be the other positions.
One term of the separation identity, with its scalar coefficient removed,
has the form
\[
 \psi_a(y_a)\psi_b(y_b)Q(t),
\]
where the three atoms are fixed on $J$, have supremum norm at most one,
and have complexity at most $H$.  In accordance with the chart calculus,
unrestricted fixed linear modulations do not count toward $H$.  Put
\[
 A_r=\sup_{\xi\in\R}
 \left|\int_{\R}Q(ru)\zeta(u)e(-\xi u)\,du\right|,
\]
where $\zeta$ is the fixed nonnegative cutoff used in the separation
lemma.  Enlarging fixed constants in the following bands has no effect
on the assertion.

\begin{lemma}[A scale sum for one separated term]
\label{lem:low-scale-sum}
Suppose the inputs in positions $a,b$ are admissible sequences associated
with $F_a,F_b$, with parameters $D_{\mathrm h},L$.  For $m\ge0$, let
$\mathcal I_m$ consist of intervals on which
$e^{-m-1}<A_{|I|}\le e^{-m}$.  Then
\begin{equation*}
\begin{split}
 &\sum_{I\in\mathcal I_m}|I|
 \left|H_I\bigl(\psi_a z_{a,I},\psi_b z_{b,I};Q\bigr)\right|\\
 &\qquad\le
 e^{-m}\operatorname{poly}\bigl(1+m,\log(2H),
                 \log(2D_{\mathrm h}),1+L\bigr)
 \|\1_JF_a\|_2\|\1_JF_b\|_2.
\end{split}
\tag{5}\label{eq:src-5}
\end{equation*}
The notation on the left denotes the two-input integral whose kernel is
$w_I(x,t)Q(t)$, with the displayed multipliers in their respective
positions.  All intervals are contained in $J$.  If $A_{|I|}=0$, the
corresponding two-input integral vanishes.
\end{lemma}

\begin{proof}
Set
$\mathcal K=1+m+\log(2H)$.  Throughout this proof, $D_1,D_2,\ldots$
denote quantities with logarithms bounded by a polynomial in
$\mathcal K$.  The single-scale part of
Lemma~\ref{lem:single-scale-separation} gives the bilinear bound $CA_r$ before
any expansion.  When $A_r>0$, Lemmas~\ref{lem:weyl} and
\ref{lem:chart-expansion} give an expansion into terms
\[
 e(\lambda_*t^2+\gamma t),\qquad |\lambda_*|r^2\le D_1.
\]
The possible unscaled $\lambda_*$ and $\gamma$ are fixed rational-height
evaluations of the original parameters.  In particular, they do not
change their underlying parameter vector when $r$ changes.

We first justify a finite truncation in a norm appropriate for a scale
sum.  Use the pooled tail conclusion of
Lemma~\ref{lem:chart-expansion} to make the sum of the worst-case absolute
tail coefficients, over all the effective quadratic and linear choices,
smaller than
$e^{-2m}(1+D_1)^{-2}$, with any additional fixed polynomial reserve
needed below.  For each fixed tail pair $(\lambda_*,\gamma)$, write
$e(\lambda_*t^2)=1+(e(\lambda_*t^2)-1)$.  The linear part is estimated
by the mean-zero square functions proved below, with a fixed modulation
$\gamma$.  The second part has smooth-kernel norm at most
$C|\lambda_*|r^2\operatorname{poly}(D_1)$; increasing the tail reserve
handles this polynomial.  At a fixed $\lambda_*$,
\[
 \sum_{r:\,|\lambda_*|r^2\le D_1}|\lambda_*|r^2\le C D_1.
\]
At each depth, contractions and Cauchy--Schwarz give the requisite product
of ordinary $L^2$ norms; for an energy sequence the same argument costs
at most $1+L$.  Thus the tail is acceptable in the absolute scale sum,
not merely at an individual interval.

There is now a finite main catalog, with absolute coefficient norm and
rational heights bounded by $D_2$.  Choose a sufficiently small
$\epsilon>0$ after that catalog, with $\log(1/\epsilon)$ polynomial in
$\mathcal K$.  Mark as exceptional every scale where a possible
nonzero $|\lambda_*|r^2$, or a possible nonzero
$r|\gamma-\gamma'|$, belongs to a prescribed intermediate annulus
between $\epsilon$ and $\epsilon^{-1}$.  Enlarge the annulus by fixed
powers of the catalog bounds when necessary.  The height-bin observation
from Section~\ref{sec:compression} bounds the number of exceptional
scales by a polynomial in $\mathcal K$.  At these scales use the original
single-scale bound $CA_r\le Ce^{-m}$, before summing coefficients.
This gives the desired estimate for the exceptional scales.

At each remaining scale the effective quadratic is negligible, and the
linear centers partition into clumps of diameter at most $\epsilon/r$,
separated by more than $\epsilon^{-1}/r$.  One way to arrange the stated
diameter is to use a smaller tiny threshold than $\epsilon$ by the main
catalog count, and then take connected components of the tiny-distance
relation.  The enlarged intermediate annulus excludes all intermediate
distances.  Replace every effective quadratic by zero and every center
in a clump by one representative.  The height-bin observation gives
\[
 \sum_r\max\{r^2|\lambda_*|:\ r^2|\lambda_*|\text{ is tiny at }r\}
 +\sum_r\max\{r|\gamma-\gamma'|:\
                       r|\gamma-\gamma'|\text{ is tiny at }r\}
 \le \epsilon\operatorname{poly}(\mathcal K),
\]
after readjusting $\epsilon$ by catalog factors.  Multiplication by the
absolute coefficient norm therefore makes the total replacement error
smaller than the right side of \eqref{eq:src-5}.

Every combined clump coefficient is $O(e^{-m})$.  Indeed, test the
original expansion of $Q(r\,\cdot)$ against
$\zeta(u)e(-r\gamma_{\mathcal C}u)$, where
$\gamma_{\mathcal C}$ is the chosen representative.  The coefficient
of its own clump is multiplied by the nonzero constant $\int\zeta$.
All other clumps give rapidly decreasing Fourier tails, and the
quadratic, truncation, and replacement errors have already been made
smaller than $e^{-m}$ by the choice of $\epsilon$ and the tail reserve.
The defining Fourier supremum is at most $e^{-m}$.

For completeness, we prove the square-function estimate used for these
linear clumps.  Make the fixed change of coordinates
$(x,t)\mapsto(y_a,y_b)$.  In these coordinates, $w_I$ has zero integral
in either variable while the other is fixed.  Its double primitive has
compact support in an interior rectangle of $I\times I$.  Expand that
primitive in a smooth Fourier series on a fixed larger rectangle,
multiply by interior cutoffs equal to one on its support, and
differentiate once in each variable.  This expresses the kernel as a
sum of tensor products of smooth, interior, mean-zero bumps.  Their
coefficients decrease faster than every fixed polynomial in the tensor
indices, while their fixed-order derivative bounds grow at most
polynomially.  It is consequently enough to consider rows
\[
 \mathcal W_{I,\nu}F
   =|I|^{-1/2}\int_I e(\nu y)\omega_I(y)F(y)\,dy,
 \qquad \nu=\pm\gamma/(b-a),
\]
where $\omega_I$ is such a bump.  The fixed bounded multiplier $\psi_a$
or $\psi_b$ can be incorporated into each row.

These rows are Bessel with squared bound polynomial in $\mathcal K$,
jointly over the retained scales, intervals, and clump representatives.
Here are explicit frequency estimates.  Decompose the input for a row
into the demodulated shells
$|\xi+\nu|\asymp2^\ell/|I|$.  At fixed scale and center, mean zero and
the $L^2$ Poincar\'e inequality give row-operator norm $O(2^\ell)$ when
$\ell<0$.  When $\ell\ge0$, two frequency antiderivatives and
integration by parts give $O(2^{-2\ell})$.  These estimates are summed
over the disjoint intervals at that scale.
For a fixed Fourier variable $\xi$ and shell index $\ell$, the
height-bin observation, applied to the original horizontal parameters,
the fixed modulation, and $\xi$, bounds the number of possible scales
by $\operatorname{poly}(\mathcal K)$.  At one scale, separation of the
centers bounds their shell multiplicity by $O(1+2^\ell)$.  Thus the
total frequency multiplicity is
$\operatorname{poly}(\mathcal K)(1+2^\ell)$.  Taking its square root
and summing the preceding row norms over $\ell$ proves the Bessel
estimate.  Exact center frequencies form a null set and cause no
additional term.  A fixed bounded multiplier preserves the estimate.

For a projection sequence, apply Lemma~\ref{lem:compression} to these
rows.  Its cost is polynomial in $\log(2D_{\mathrm h})$.  For an energy
sequence, use instead the per-interval Bessel estimate over separated
centers and then sum its given energy bound; no composition lemma is
needed.  Cauchy--Schwarz over the two sets of rows, together with the
$O(e^{-m})$ bound for each combined coefficient, proves
\eqref{eq:src-5}.  If $A_r=0$, the cutoff Fourier transform is zero,
so $Q$ vanishes on the part of the $t$-axis relevant to the kernel.
\end{proof}

\begin{corollary}[A finite predictor list]
\label{cor:predictor-list}
Fix the two low atoms on $J$, with complexity at most $H_{\rm low}$.
For every $B\ge1$, their product in the form can be replaced, with
absolute length-weighted scale-sum error at most
\[
 e^{-B}\operatorname{poly}\bigl(\log(2D_{\mathrm h}),1+L\bigr)
 \|\1_JF_a\|_2\|\1_JF_b\|_2,
\]
by a finite list of terms
\begin{equation*}
 \psi_a(y_a)\psi_b(y_b)w_I(x,t)e(\lambda_*t^2),
 \qquad |\lambda_*||I|^2\le D_{\rm pred}.
\tag{6}\label{eq:src-6}
\end{equation*}
The scalar coefficients may depend on the scale.  The list of triples
$(\psi_a,\psi_b,\lambda_*)$ is fixed on $J$.  Its size, the absolute
coefficient bounds, and all atom complexities satisfy
\[
 \log D_{\rm pred}
 \le\operatorname{poly}\bigl(1+B,\log(2H_{\rm low})\bigr).
\]
The multipliers $\psi_a,\psi_b$ may be normalized to have supremum norm
at most one by incorporating their sizes into the coefficients.
\end{corollary}

\begin{proof}
Expand the low pair by Lemma~\ref{lem:ap-separation}.  Its dyadic shell
coefficients have the stated rapidly decreasing absolute sums.  Apply
Lemma~\ref{lem:low-scale-sum} to each shell and sum its Fourier-supremum
bands.  First discard sufficiently high separation shells; then, for
each retained term, discard sufficiently small $A_r$ bands.  The shell
and band thresholds have logarithms polynomial in
$1+B+\log(2H_{\rm low})$.  The remaining scales admit the finite
chart expansion from Lemma~\ref{lem:chart-expansion}.  Truncate its
pooled tails in the absolute scale-sum norm exactly as in the proof of
Lemma~\ref{lem:low-scale-sum}.  Distribute the fixed modulation
$e(\gamma t)$ between positions $a,b$.  The surviving rational-height
choices of quadratic parts and linear shifts form a fixed finite list.
All of its bounds have the asserted logarithmic size.  The discarded
errors have the required sum by \eqref{eq:src-5}.
\end{proof}

\subsection{Predictor moments on a finer partition}

A predictor for the multiplier $\psi_a$ is the augmented column
\[
 \Phi_p=(\overline{\psi_a},\sqrt\epsilon\,\mathbf e_p),
 \qquad 0<\epsilon\le1.
\]
The conjugation is chosen so that the physical residual moment is
bilinear against $\psi_a$.

\begin{lemma}[Delayed predictor replacement]
\label{lem:predictor-replacement}
Consider one fixed term of \eqref{eq:src-6}, born on $J$.  Suppose the
nested predictor family in position $a$ contains $\Phi_p$ on every
use interval and on all its descendants.  Other columns are allowed,
subject to the projection budget $D_{\mathrm h}$.  Let the other
position $b$ contain an admissible sequence.  For a use interval $I$
of depth $s$, replace $F_a$ by the projection of $F_a$ on the partition
at depth $s+K_0$, where $K_0\ge0$ is an integer.  The absolute
length-weighted error, summed over any collection of uses, is at most
\begin{equation*}
 C\bigl(2^{-K_0}+(1+K_0)\sqrt\epsilon\bigr)
 \operatorname{poly}\bigl(D_{\rm pred},\log(2D_{\mathrm h}),1+L\bigr)
 \|\1_JF_a\|_2\|\1_JF_b\|_2.
\tag{7}\label{eq:src-7}
\end{equation*}
The assertion also permits sequential replacement in positions $a,b$.
\end{lemma}

\begin{proof}
The predictor identity expresses the multiplier-weighted mean of the
physical residual as $\epsilon$ times the predictor's column
coefficient.  Smooth-bump approximation below the delayed partition
supplies the $2^{-K_0}$ term, while subtracting the coarse mean and
summing private-coordinate increments gives the
$(1+K_0)\sqrt\epsilon$ term.

Let $E_h$ denote conditional expectation on the depth-$h$ partition
of $J$, and write $P_h$ for the direct sum of its local projections.
On descendants where the predictor is present, put
\[
 D_h=(f_a-(P_hF_a)_{\phys})\psi_a,
 \qquad \mu_h=E_hD_h.
\]
If $a_{h,p}$ is the predictor coefficient on a depth-$h$ cell,
orthogonality to $\Phi_p$ gives
\[
 \mu_h=\epsilon a_{h,p}.
\]
Consequently $\mu_{h+1}-\mu_h$ is $\sqrt\epsilon$ times the change
of the $p$-th private coordinate of $P_hF_a$.  This identity continues
to hold after arbitrary further column insertions.

For a smooth normalized bump $\omega_I$ on an interval of depth $s$,
set $H=s+K_0$.  The residual moment has the exact telescoping
representation
\[
\begin{split}
 \int_I D_H\omega_I
  ={}&\int_I\mu_H\omega_I\\
   &+\sum_{h\ge H}\int_I(D_h-D_{h+1})(\omega_I-E_h\omega_I)\\
   &+\sum_{h\ge H}\int_I(\mu_{h+1}-\mu_h)
                                  (E_{h+1}\omega_I-E_h\omega_I).
\end{split}
\]
To verify it, telescope
$\int_I D_h(\omega_I-E_h\omega_I)$ and use
$E_{h+1}D_{h+1}=\mu_{h+1}$; the term involving $\mu_h$ and
$E_{h+1}\omega_I-E_h\omega_I$ vanishes by conditional expectation.
The final remainder tends to zero because $D_h$ has uniformly bounded
$L^2(I)$ norm and the smooth bump's conditional-expectation error tends
to zero.  Thus no convergence of $P_hF_a$ to $F_a$ is being assumed.

The smoothness of the bump gives a factor $O(2^{s-h})$ in each of
the two infinite sums.  Its other factor is a localized orthogonal
projection increment, with an additional $\sqrt\epsilon$ for the
private-coordinate term.  At a fixed depth offset $h-s$, sum squared
row norms over intervals of depth $s$, then sum over $s$ using
orthogonality of the increments $P_{h+1}-P_h$.  Minkowski's inequality
over the offsets gives an aggregate row norm
\[
 C2^{-K_0}\|\1_JF_a\|_2.
\]
This estimate uses only the bounded multiplier $\psi_a$, and is
unaffected by the size of the enlarged column catalog.

If $\omega_I$ has mean zero, the first term can be written with
$\mu_H-\mu_s$.  Decomposing it into the $K_0$ successive private
increments gives aggregate row norm at most
\[
 C K_0\sqrt\epsilon\,\|\1_JF_a\|_2.
\]
Apply the mean-zero tensor decomposition from the preceding proof
to $w_I$, including the fixed distributed modulation in the fixed
multipliers.  Pair the resulting error rows in position $a$ with
the Bessel rows in position $b$.  Compression handles a current
projection in that position; an energy sequence is handled by its
per-interval Bessel estimate.  This proves
\eqref{eq:src-7} for the part of the kernel with
$e(\lambda_*t^2)$ replaced by one.

For the difference kernel
$w_I(e(\lambda_*t^2)-1)$, an ordinary smooth tensor expansion suffices.
Its derivative costs contain the factor
$|\lambda_*|r^2\operatorname{poly}(D_{\rm pred})$.
The same telescoping argument estimates its expectation-free part.
For the initial $\mu_H$ term, which need not be tested against a
mean-zero bump, use the per-depth bound
$C\sqrt\epsilon\|\1_JF_a\|_2$ and sum
$|\lambda_*|r^2$ geometrically over the scales satisfying
$|\lambda_*|r^2\le D_{\rm pred}$.  This accounts for the additional
$\sqrt\epsilon$ term in \eqref{eq:src-7}.

For sequential replacement, write a previously delayed projection
as its current projection plus a rolling projection difference.
The latter has a square-energy bound with $L\le C(1+K_0)$, by
telescoping its orthogonal increments.  It is therefore an admissible
opponent in the second application of the argument.
\end{proof}

\subsection{Adaptive approximants at staggered accuracies}

For a physical residual $r$ in position $j$, let
$\mathfrak d_{j,I}(r)$ be the supremum of $|H_I|$ when that position
is occupied by $r$ and the other positions contain one normalized
$L^q(I)$ input and two normalized $L^2(I)$ inputs.  The supremum
also ranges over which opposing position receives the $L^q$ input.

\begin{lemma}[Inherited adaptive approximation]
\label{lem:adaptive-approximation}
Suppose $\|f_j\|_{q,I}\le1$ throughout a finite normalized subtree.
For every sufficiently large $k$, one can construct, on segments
whose roots have bounded total length, local projections $g_I=P_IF_j$
such that
\[
 \mathfrak d_{j,I}\bigl(f_j-(g_I)_{\phys}\bigr)\le e^{-k},
 \qquad \|g_I\|_{2,I}\le1.
\]
Within each segment the columns are inherited and the projections
are nested in depth.  Column path counts, atom complexities, and
the physical supremum bounds of $g_I$ are at most
$\exp(k^{O(1)})$.  Fresh columns may all have private weight one.
The polynomial degrees are independent of $q$.
\end{lemma}

\begin{proof}
Start with no columns at a segment root $J$ and proceed in depth
order.  At a cell whose residual fails the desired inequality,
Lemma~\ref{lem:detection} supplies a bounded atom $\varphi$ with
physical residual correlation at least $\rho$, where
$\log(1/\rho)\le k^{O(1)}$ and the atom complexity obeys the same
logarithmic bound.  Adjoin
$(\varphi,\mathbf e_v)$ with a fresh private coordinate and refit.
The old residual has zero component in this fresh direction.
The squared projection energy therefore increases by at least
$c\rho^2|I|$.  Repeated fits at one cell and refining fits at successive
depths form one nested sequence of subspaces in $L^2(J)$, so their
total energy gain is at most $\|\1_JF_j\|_2^2\le|J|$.

Stop at the first cells where the cumulative number of insertions
along a path would exceed $C\rho^{-2}$.  Integrating the insertion
count over the disjoint first-excess cells bounds their total length
by $C^{-1}$ times a fixed multiple of $|J|$.  Choose $C$ so this is
at most $|J|/4$, and restart the same procedure on those roots.
Iteration gives a geometric bound for the sum of segment-root
lengths.  On every retained cell the greedy fit has terminated with
the required residual estimate.

With private weight one, the coefficient $\ell^2$ norm is bounded
by the normalized projection norm.  A catalog of at most
$\exp(k^{O(1)})$ bounded scalar columns consequently gives the
claimed physical supremum bound by Cauchy--Schwarz.  Complexity and
count bounds follow from detection and the insertion cap.  The
uniform polynomial-degree assertion is inherited from
Lemma~\ref{lem:detection}.
\end{proof}

Use levels $t=1,2,\ldots$, updating positions in a fixed cyclic
order.  Choose an absolute $A>1$ and large initial $k_1$, and set
$k_{t+1}=k_t^A$.  Apply Lemma~\ref{lem:adaptive-approximation}
at level $t$ with accuracy $e^{-k_t}$; denote its projection in
the updated position by $g_{t,I}$.  Put $g_{t-4,I}=0$ when the
position has not yet been updated.  The columns for a fixed level
are inherited except at the sparse resets just constructed.
When finitely many levels are under discussion, intersect their
segments.  For a fixed bounded number of levels this common
refinement still has bounded total root length: every new common
root is a root of one of the constituent segment families, or the
original root.

\begin{lemma}[The two-earliest-increments expansion]
\label{lem:two-accuracy-expansion}
Take the two earliest increment levels $u<t$ in distinct positions
$i,j$, and let $a,b$ be the complementary positions.  Define
\begin{equation*}
 X_i=g_u-g_{u-4},\qquad X_j=g_t-g_{t-4},\qquad
 X_a=F_a-g_{a,\mathrm{prev}},\qquad
 X_b=F_b-g_{b,\mathrm{prev}}.
\tag{8}\label{eq:src-8}
\end{equation*}
Here $g_{e,\mathrm{prev}}$ is the last update in position $e$ before
level $t$, and is zero if no such update exists.  For a finite
collection of intervals, the original form is the limit of the
sum of the corresponding tuple forms $H_I(X)$ as the update cutoff
tends to infinity.  If $A$ is sufficiently large, there is an
absolute $\theta_0>0$ such that, whenever $t\ge12$ and $k=k_t$,
every pair of full arguments in \eqref{eq:src-8} has bilinear form
norm at most $e^{-2k^{\theta_0}}$ against two normalized $L^2(I)$
inputs.  All four arguments have bounded normalized $L^2(I)$ size.
\end{lemma}

\begin{proof}
For a finite terminal level, decompose each actual input into its
increments up to that level and its terminal residual.  In every
term having at least two increments, record the earliest two levels
$u<t$.  Summing all later choices in either other position gives
exactly the input minus its last update before $t$.  This yields
\eqref{eq:src-8} for all such terms.  The terms with at most one
increment consist of either four terminal residuals, or one whole
terminal projection and three terminal residuals.

We record the comparison that controls these remainders and the
claimed pair estimates.  If $h>\ell$ are updates in distinct
positions, replacing their pair of projections by the corresponding
actual inputs has form-norm error at most
\[
 C_q\bigl(e^{-k_h}\exp(C_q k_\ell^{d})+e^{-k_\ell}\bigr)
\]
for an absolute $d$.  Replace the more accurate projection first,
using the physical supremum bound of the other as its opposing
$L^q$ norm, and then replace the less accurate projection against
the actual $L^q$ input.  Choose $A>d$, with fixed reserve.  The
same reasoning applies when a projection is already replaced by
an actual input.  In particular, a pair of terminal residuals has
form norm tending to zero.  Their two remaining opponents have
bounded $L^2$ size, proving the assertion about the terminal
remainder on a finite interval collection.

For the tuple estimate, first suppose $u<t-4$.  Whenever the early
argument in position $i$ is paired with a later difference or
residual, compare that later argument directly with its actual
input against each of the two earlier projections.  The actual
terms cancel in the later difference or residual.  Its smallest
accuracy index is at least $t-4$, and exceeds the relevant earlier
indices; the stagger therefore defeats their supremum bounds.
Pairs not involving this early position use only indices at least
$t-4$, apart from the preceding projection in position $j$, whose
index is exactly $t-4$.

If instead $u\ge t-4$, distinct positions force $u\ge t-3$;
all nonzero projections in the tuple then have indices at least
$t-7$.  Expand a selected pair into projection/actual pairs and
use the preceding comparison.  At least one selected full
argument has coefficient sum zero after replacing projections by
their actual inputs, so the two-actual terms cancel.  The same
argument covers the pairs in the first case that do not contain
the early position.  Each error is bounded by a fixed multiple
of $\exp(-c_q k_{t-7})$ for large starting log.  Since
$k_{t-7}=k_t^{A^{-7}}$, decreasing a fixed exponent below
$A^{-7}$ gives the stated $e^{-2k^{\theta_0}}$ bound.  The
$L^2$ size assertion follows from contraction and the triangle
inequality.
\end{proof}

\subsection{Delayed high fits and their coarse stopping cells}

Fix one tuple and a common segment for its older lists.  The two
positions $i,j$ in \eqref{eq:src-8} are called low only to specify
which arguments will be expanded into their original columns.
The later construction does not expand all four arguments
coefficientwise.

\begin{proposition}[Truncating the two high positions]
\label{prop:high-truncation}
For the tuple above, put $k=k_t$.  There are an integer
$K_0\le k^{O(1)}$, further bounded-packing segments, and nested
projection families $P^{\rm new}_a,P^{\rm new}_b$ on those
segments, for which
\begin{equation*}
 Y_i=X_i,\qquad Y_j=X_j,\qquad
 Y_e=P^{\rm new}_{e,s+K_0}F_e-g_{e,\mathrm{prev}}\quad(e=a,b),
 \qquad s=\operatorname{depth}(I).
\tag{9}\label{eq:src-9}
\end{equation*}
The absolute length-weighted replacement error in $H_I$, summed
over the original normalized root, is at most $Ce^{-k}$ times
its length.  Path counts, inverse private weights, atom
complexities, and the physical supremum bounds on coarse
evaluation cells are all at most $\exp(k^{O(1)})$.
For $t\ge12$, after decreasing an absolute $\theta>0$ if necessary,
every pair of full $Y$ arguments has form norm at most
$\eta=e^{-k^\theta}$ against normalized $L^2$ opponents.
All $Y$ arguments have bounded normalized $L^2$ size.
\end{proposition}

\begin{proof}
We give the choices in their dependency order.  On a common older
segment, the low columns, their coefficients, and their path
counts have a bound $D_{\rm old}=\exp(k^{O(1)})$.  Apply
Corollary~\ref{cor:predictor-list} to every available pair of low
labels, using a precision log $B$ that is a sufficiently large
fixed polynomial in $k$.  Insert the conjugate predictor columns
in positions $a,b$ at the first joint availability of that pair,
and retain them thereafter.  Pairs inherited at the segment root
have their predictors present from that root.  Use private weight
$\sqrt\epsilon$ for these predictors.

Choose $K_0$ and $\log(1/\epsilon)$ to be sufficiently large
fixed polynomials in $k$, after $D_{\rm pred}$ and the original
low coefficient/count bounds have been fixed.  The choices make
both the error of Corollary~\ref{cor:predictor-list} and the
error of Lemma~\ref{lem:predictor-replacement} smaller than
$e^{-2k}$ after all original-label sums.  Only a polynomial in
the logarithm of a subsequently enlarged high catalog occurs in
these errors.  Thus the choices remain valid under any later
fixed-polynomial growth of that log catalog bound, by increasing
the starting $k$.

Here is the summation over low labels in detail.  Expand the two
low projections into their original labels only in this
application.  A fixed pair is available on the subtree rooted
at its deeper birth cell, or on the segment root if it was
inherited.  Its coefficients on a use cell are constant and
bounded by a fixed power of $D_{\rm old}$.  The number of such
pair-birth roots on a path is also bounded by a fixed power of
$D_{\rm old}$.  Consequently, on a normalized segment root $R$,
\[
 \sum_{J\text{ pair-birth root}}
        \|\1_JF_a\|_2\|\1_JF_b\|_2
 \le D_{\rm old}^{C}|R|.
\]
Indeed Cauchy--Schwarz bounds this sum by the product of the
square roots of the corresponding two sums of squared masses,
and the path multiplicity bounds each by
$D_{\rm old}^C\|\1_RF_e\|_2^2$.  Apply
Corollary~\ref{cor:predictor-list} on both sides of each
replacement, and then Lemma~\ref{lem:predictor-replacement}
to each retained term.  A previously delayed opposing
projection is its current projection plus a rolling
square-energy sequence with $L=O(1+K_0)$.  This proves the
absolute-sum error assertion, even after the extra fits now
described.

The predictors alone need not give the pair estimates in the
statement.  We therefore add tuple-specific accurate fits.
Proceed chronologically in projection depth $h$.  At depth
$h$, first admit any predetermined predictor columns becoming
available there.  A coarse decision cell at that step has
depth $s=h-K_0$.  Test the physical residual
\[
 f_e-(P^{\rm new}_{e,h}F_e)_{\phys}
\]
on that coarse cell in the norm $\mathfrak d_{e,I}$.
The normalized $L^2(I)$ size of this residual is bounded,
because $P^{\rm new}_{e,h}$ is a contraction on the partition
of $I$ into its depth-$h$ descendants.  Require a precision
whose log is a sufficiently large fixed power of $k$.
On failure, detection gives a bounded scalar column with a
coarse residual moment at least $\rho'$, where
$\log(1/\rho')\le k^{O(1)}$.  Insert one new common label
with private weight one on all the depth-$h$ subcells of
$I$, and retain that label on their descendants.

The joint gain from this insertion is at least
$c(\rho')^2|I|$.  To see this, let $m_Q$ be the normalized
residual moment on a depth-$h$ subcell $Q\subset I$.
The fresh private direction is absent from every old tested
residual, so the gain on $Q$ is at least $c|Q||m_Q|^2$.
Cauchy--Schwarz and the coarse moment bound give
\[
 \sum_{Q\subset I,\,\operatorname{depth}(Q)=h}
                  c|Q||m_Q|^2
 \ge c|I|\left|\fint_I\text{residual}\,\overline{\varphi}\right|^2
 \ge c(\rho')^2|I|.
\]
All ancestor extra requests affecting a coarse cell are common
to its fine partition.  Cap their cumulative number along a
coarse path at a sufficiently large multiple of
$(\rho')^{-2}$.  First excess coarse cells have small total
relative length, by the joint gains and the orthogonal energy
bound for the entire chronological iteration.  Send their
coarse summands and descendants to fresh attempts.  Predictor
admissions have their separate path cap from the older lists
and the finite predictor catalogs.

There is no conflict between a coarse cut and a delayed
projection already used above it.  If a coarse cell is cut
when projection depth is $h$, all its strict ancestors have
already had their delayed evaluations computed, since their
projection depths are smaller than $h$.  Continue the stopped
estimator on the cut subtree by inheritance and refitting,
with its predictor moments retained; its later residual tests
are no longer needed for the summands sent to the new
attempt.  This extension supplies all finer projections
needed in the telescoping proof of
Lemma~\ref{lem:predictor-replacement}.  A fresh attempt for
the descendant coarse summands may use a new extra-label
catalog.  Intersect the resulting coarse segment families
between the two high positions and the older lists.
Their roots still have bounded total length, and each
resulting piece uses restrictions of genuine nested
projection sequences, including their extensions beyond
its boundaries.

The accuracy order is as follows.  Choose the first new
high precision to dominate all previous projection and
predictor costs.  Detection fixes the first new catalog
bound.  Choose the second new high precision to dominate
that first catalog's physical supremum bound as well.
For a fine subcell $Q\subset I$ at depth $s+K_0$,
\[
 \|F_e\|_{2,Q}\le 2^{K_0/2}\|F_e\|_{2,I}.
\]
The private-coordinate coefficient estimate thus bounds
the first new projection's physical supremum on $I$ by
$2^{K_0}$ times a fixed power of its own count and
conditioning bound.  It does not involve the next high
precision.  The two choices are therefore successive
fixed-polynomial log choices.

Apply the comparison argument of
Lemma~\ref{lem:two-accuracy-expansion} to these two new
approximants and the older projections.  The new fits are
more accurate than every older projection, and the second
is more accurate than the first with sufficient reserve
for its supremum bound.  In each pair of full arguments
in \eqref{eq:src-9}, the actual-input contributions cancel
exactly as before.  If the earliest low position has a
much earlier index, the other selected full argument is
a difference of later approximants and is tested directly
against that earlier position.  Otherwise all older
indices involved are at least $t-7$.  Decreasing the
absolute exponent from
Lemma~\ref{lem:two-accuracy-expansion} yields the claimed
$\eta$ bound.  Contraction on the fine partition and on
the older coarse cells gives the local $L^2$ bounds.
Finally, the new high fits are specific to this tuple;
their larger precisions do not enter the global update
recurrence $k_{t+1}=k_t^A$.
\end{proof}

\subsection{Base sides and the active collection}

Work on one normalized segment root $R$ from
Proposition~\ref{prop:high-truncation}.  Set
\[
 Y_l=B_l+U_l,\qquad U_i=U_j=0,
 \qquad
 U_e=(P^{\rm new}_{e,s+K_0}-P^{\rm new}_{e,s})F_e
 \quad(e=a,b).
\]
Thus each $B_l$ is a difference of at most two current
basic projections.  These basic columns are inherited
within the segment, and labels in different basic
catalogs have distinct private directions.  Choose
$D_{\rm base}$, with
\[
 e^k\le D_{\rm base}\le\exp(k^C),
\]
to dominate the base path counts, inverse private
weights, complexities, and current coefficient sums.
The local normalized sizes of $B,U,Y$ are bounded.
Rolling orthogonal projection differences give, also
on any normalized subroot,
\[
 \sum_{I\subseteq R}|I|\|U_{e,I}\|_{2,I}^2
 \le C(1+K_0)^2|R|.
\]
For example, expand the rolling difference into its
$K_0$ orthogonal increments and sum at fixed depth
offset; the triangle inequality in the resulting
row space gives the displayed bound.

The term containing both possible $U$ arguments is
already summable.  Its complementary two arguments
are the full $Y_i,Y_j$, whose pair norm is at most
$\eta$.  Cauchy--Schwarz over intervals therefore
bounds its absolute length-weighted sum by
$C\eta(1+K_0)^2|R|$.  For the finitely many leading
levels the same assertion holds with the factor
$\eta$ omitted.

\begin{lemma}[Few active scales for a fixed triple]
\label{lem:active-triple}
Fix three bounded basic scalar columns, in three
distinct positions, on their common birth subtree.
For a dyadic interval $I$ let $\mathcal A_I$ be the
norm of their form against a normalized $L^2(I)$
input in the fourth position.  Along every spatial
path, the number of scales for which
$\mathcal A_I>\delta$ is polynomial in
$\log D_{\rm base}+\log(1/\delta)$.
Moreover, for $0<\delta\le1$,
\[
 \sum_{I\text{ on the path}:\,\mathcal A_I\le\delta}
          \mathcal A_I
 \le \delta\operatorname{poly}
              \bigl(\log D_{\rm base},\log(1/\delta)\bigr).
\]
The bound is uniform in the locations of the path
and in all unrestricted oscillation parameters.
\end{lemma}

\begin{proof}
Fix a path point $y$, and let the fourth position
be $l$.  Writing the transpose at $y+h$, its
squared $L^2$ norm is an integral in the scaled
variables $(h,t,t')/r$, $r=|I|$, containing the
fixed triple and its conjugate copy and the two
smooth kernels at the shared anchor.  The six
atoms combine into the multivariable chart
setting of Lemma~\ref{lem:weyl}.  Recenter at
the fixed $y$.  All unscaled horizontal, linear,
and quadratic parameters are now fixed along
the path, before multiplication by $r$ or $r^2$.

If $\mathcal A_I>\delta$, apply
Lemma~\ref{lem:weyl} to this integral at Fourier
frequency zero, and then
Lemma~\ref{lem:chart-expansion} with an error
small compared with $\delta^2$.  Every effective
quadratic entry and every linear $t$-frequency
in the truncated expansion is of the form
$r^2 a\cdot\chi_*$ or $r a\cdot\chi_*$,
respectively, for a fixed finite parameter vector
$\chi_*$.  The coefficient vectors have rational
height quasipolynomial in $D_{\rm base}$ and
$\delta^{-1}$; the dimension is polynomial in
the original chart dimension.  This description
also covers the choices of rational splitting.

Choose tiny/large gap thresholds after the finite
truncation, with reserve for its absolute
coefficient norm.  Outside polynomially many
intermediate scales, the height-bin observation
makes every effective quadratic entry that is
bounded in the Weyl expansion tiny, and every
retained $t$-frequency either tiny or large.
Remove the effective quadratic by uniform
approximation.  Large $t$-frequencies integrate
to a negligible amount by smoothness of the
kernel.  For a tiny $t$-frequency, replace its
oscillation by a constant; the remaining
$t$-integral is zero by the exact marginal
\[
 \int w_I(y+h-lt,t)\,dt=0
\]
at fixed $h,t'$.  The other kernel and the
remaining linear oscillations are independent
of this $t$ variable.  All errors, after
coefficient summation, are smaller than
$\delta^2$.  This contradicts
$\mathcal A_I>\delta$ and proves the count.

Apply that count to the successive bands
$2^{-n-1}\delta<\mathcal A_I\le2^{-n}\delta$,
using threshold $2^{-n-1}\delta$.  Their total
contribution is bounded by
\[
 \sum_{n\ge0}2^{-n}\delta\,
 \operatorname{poly}\bigl(\log D_{\rm base}
                          +\log(1/\delta)+n\bigr),
\]
which is the asserted small-norm sum.  The
kernel may vary with $I$ within its fixed class:
only its uniform smoothness and exact marginal
were used.
\end{proof}

\begin{proposition}[Reduction to the active cell minimum]
\label{prop:base-reduction}
The remaining patterns, with at least three $B$
arguments and optionally one $U$ argument, can
be restricted, with total absolute error at
most $e^{-k}|R|$, to a common active collection
whose path multiplicity satisfies
\begin{equation*}
 M\le D_{\rm base}^{C}.
\tag{10}\label{eq:src-10}
\end{equation*}
The finitely many leading update levels have
a bound independent of the number of depths.
For $t\ge12$, it suffices on each segment and
for each retained pattern $z$ to estimate
\begin{equation*}
 \sum_{I\ \mathrm{active}}|I|
                       \min\bigl(\eta,|H_I(z)|\bigr).
\tag{11}\label{eq:src-11}
\end{equation*}
The two nonside arguments complementary to a
pair of later side arguments are the whole
$B$ arguments of this tuple before further
filters.  Restoring them to their full $Y$
arguments recovers the pair-smallness above.
Further attempts may restart on normalized
subroots with these basic arguments and their
inherited lists unchanged.
\end{proposition}

\begin{proof}
Declare $I$ active if some triple of available
pure basic scalar columns in three positions
has transpose norm greater than
$D_{\rm base}^{-C'}$, with $C'$ a sufficiently
large fixed constant.  Along a path there
are at most $D_{\rm base}^3$ triples, and
Lemma~\ref{lem:active-triple} gives only a
logarithmic-polynomial number of active
scales for each one.  Enlarging the absolute
power proves \eqref{eq:src-10}.

On an inactive interval, expand three of
the $B$ arguments into their basic columns.
Their coefficient sums and the local $L^2$
size of the remaining argument cost only
a fixed power of $D_{\rm base}$.  For a
fixed triple, the small-norm assertion of
Lemma~\ref{lem:active-triple} bounds its
path sum by
$D_{\rm base}^{-C'}
 \operatorname{poly}(\log D_{\rm base},C')$.
Sum over the pathwise triple count and
integrate over path points.  Choosing
$C'$ larger than all these fixed power
costs, and using $D_{\rm base}\ge e^k$,
gives the stated $e^{-k}|R|$ error.

For each fixed leading level, all budgets
and the multiplicity in \eqref{eq:src-10}
are fixed constants.  Local absolute
$L^2$ estimates and bounded sizes thus
give a bound independent of the number
of depths.  For later levels, the pair
estimate in
Proposition~\ref{prop:high-truncation}
and the bounded sizes of the other two
arguments give $|H_I(Y)|\le C\eta$.
Decrease $\theta$ once more, and increase
the starting log, to absorb this fixed
constant into the definition of $\eta$.
Multilinearity expands $H_I(Y)$ into
the retained patterns and the already
estimated term with both $U$ arguments.
For nonnegative quantities,
$\min(\eta,\sum_j a_j)\le
\sum_j\min(\eta,a_j)$.  The absolute
errors and the two-$U$ term may therefore
be taken outside the minimum, leaving
\eqref{eq:src-11} separately for the
remaining patterns.

This reduction retains the whole coarse
arguments $B_l$.  In particular, later
filter removal in two nonside positions
returns exactly these arguments;
substituting $Y_l=B_l+U_l$ changes the
term only by terms with an additional
basic side.  The full complementary
pair then has the pair-smallness already
proved.  All relevant current families
remain inherited under restriction to
a normalized subroot, and their rolling
square estimates follow anew from
orthogonal energy there.  Hence a later
stopping attempt can use the same base
data.
\end{proof}

A bound of the form
$\eta^c(1+C_N)|R|$, for some fixed $c>0$,
for \eqref{eq:src-11} is sufficient for
the normalized absorption argument.
The original-label replacements have
absolutely summable errors, segment
roots pack, and the number of earlier
indices $u<t$ is at most $t$.  The
completion of the proof will sum the
resulting positive powers of
$e^{-k_t^\theta}$.
